\subsection{Temporal extraction, soundness composition, and remaining real-world gaps}\label{sec:realize:security}

\begin{definition}[Extraction-before-challenge contract]\label{def:realize:early-extract}
A concrete execution satisfies the
\emph{early extraction contract} if there is a
computational coupling with an ideal run
in which, except with probability at most
$\varepsilon_{\rm ext}$:
\begin{enumerate}[label=\textup{(\roman*)}]
\item every bit commitment whose value is
used in the first-challenge discrepancy is
assigned a unique extracted binary opening
\emph{before} the verifier samples $\theta$;
\item every auxiliary bit commitment whose value
enters the anchored residual is assigned a
unique extracted binary opening
\emph{before} the verifier samples $z$;
\item every accepted compiled parity OR proof
has an extracted aggregate opening
$(t,R)$ with $t$ in the public
permitted candidate set; every
bitness proof has an extracted binary
opening; and every packed-column witness
used by the outer protocol is linked
to its registered binary coordinates
by the separately checked packing relation;
\item the resulting coupled ideal protocol
has the same public messages and the same
fresh independent $\theta,z$ as the
real execution outside the bad event.
\end{enumerate}
The contract is a \emph{knowledge and
timing property to be proved for an actual
composition of the OR extractors},
not a consequence of bare computational
binding alone.
\end{definition}

\begin{lemma}[Accepted parity equations are valid under extracted consistency]\label{lem:realize:extracted-correct}
Suppose all bit and aggregate OR proofs have
successful extracted openings and the adversary
has not solved the setup discrete logarithm.
Then each accepted compiled affine equation
holds for the extracted registered field values.
\end{lemma}
\begin{proof}
If one accepted affine equation failed,
Theorem~\ref{thm:realize:compiler-extract}
would turn the extracted openings for its
nonzero parity row into a solution of the
setup discrete-logarithm problem.
This contradicts the stipulated
non-solution event. Hence every such
equation holds.
\end{proof}

\begin{theorem}[Conditional soundness of the explicit bit-commitment compiler]\label{thm:realize:outer-sound}
Assume the false public evaluation
instance and timing hypotheses of
Theorem~\ref{thm:zkrs:false-sound}.
Suppose the concrete bit-commitment
and OR-proof implementation satisfies
the early extraction contract of
Definition~\ref{def:realize:early-extract}
with bad-event probability at most
$\varepsilon_{\rm ext}$, and the chance
of breaking the discrete-logarithm
binding assumption in this coupling is
at most $\varepsilon_{\rm DL}$.
Then
\[
 \boxed{\Pr[\mathrm{accept\ false\ claim}]
 \leq\min\left\{1,\frac{\ell+n}{q}
                   +\varepsilon_{\rm ext}
                   +\varepsilon_{\rm DL}\right\}.}
\]
\end{theorem}
\begin{proof}
Remove the event that the early
extraction contract fails and the
event that a forbidden
discrete-logarithm opening collision
is produced. Their total probability
is at most
$\varepsilon_{\rm ext}+\varepsilon_{\rm DL}$
by the union bound.
On the complementary event,
Lemma~\ref{lem:realize:extracted-correct}
makes every accepted compiled
affine equation true on the fixed
extracted field values, with the
required order of commitments and
random challenges.
The coupled execution is therefore
a legal instance of the ideal
algebraic checker in
Theorem~\ref{thm:zkrs:false-sound}.
That theorem bounds false-claim
acceptance on this good event
by $(\ell+n)/q$.
Add the two removed bad-event
probabilities and cap by one.
\end{proof}

\begin{lemma}[Multi-statement extraction cannot be silently assumed]\label{lem:realize:extract-needs-proof}
The two-transcript special soundness
of an individual Schnorr OR proof
does not, as a matter of logical
implication alone, establish the
early extraction contract for
an arbitrary adaptive multi-stage
protocol using many proofs.
\end{lemma}
\begin{proof}
Theorem~\ref{thm:realize:or-extract}
has hypotheses concerning two
accepting transcripts of one fixed
statement sharing exactly one
first-message vector and carrying
different global challenges.
Definition~\ref{def:realize:early-extract}
instead quantifies over many
interleaved statements and their
dependency order relative to
two outer challenges and a fixed
witness representation.
The former statement contains
neither the latter's challenge
timing hypothesis nor a coupling
of the extracted witnesses across
distinct proofs. Hence the latter
requires additional premises and
a separate composition proof.
\end{proof}

\begin{definition}[Malicious sigma-verifier security target]\label{def:realize:mv-target}
The \emph{malicious-sigma-verifier
simulation target} removes the
requirement that each global
Schnorr OR challenge be independently
uniform. The verifier may compute
it as an arbitrary efficient function
of its entire view, including the
OR first messages, commitments,
outer challenges and prior proofs.
The target requires one efficient
online witness-independent simulator
for the full resulting interactive
verifier view.
\end{definition}

\begin{proposition}[The proven HVZK simulator does not discharge the malicious target]\label{prop:realize:hv-not-mv}
The proof of
Theorem~\ref{thm:realize:full-hvzk}
does not establish
Definition~\ref{def:realize:mv-target}.
\end{proposition}
\begin{proof}
The simulator of
Theorem~\ref{thm:realize:or-hvzk}
first chooses a uniform vector
of branch challenges and
defines the OR first messages
from that choice. Consequently
the global challenge it outputs
is uniformly distributed and
is not obtained by running an
arbitrary verifier on the
already determined first messages.
Definition~\ref{def:realize:mv-target}
permits the verifier to choose
its challenge by such an
arbitrary function. The
simulator's argument therefore
does not cover that quantifier.
No impossibility of an alternative
simulator is claimed.
\end{proof}

\begin{definition}[Concrete ZK compiler obligation]\label{def:realize:mv-compiler}
A \emph{malicious-verifier
zero-knowledge sigma compiler}
for this chapter is an explicit
transformation of the bitness
and parity proof transcripts
together with a proof that:
(i) its complete visible
transition kernels satisfy
Definition~\ref{def:zkrs:realization}
for all allowed malicious
verifiers; (ii) its knowledge
extraction respects the early
challenge-timing contract;
and (iii) it retains binding
to the same committed witness.
The transformation may use a
designated-verifier technique,
a proven malicious-verifier
interactive compiler, or an
appropriately modeled
Fiat--Shamir transform
\cite{FiatShamir1986}, but
neither its implementation nor
its proof is supplied here.
\end{definition}

\begin{remark}[ROM and QROM are different proof obligations]
A Fiat--Shamir transform of these
Schnorr-style proofs is not
automatically a quantum-random-oracle
proof, and conventional discrete-log
commitments are not quantum-resistant.
To claim a post-quantum KBFold
extension, the commitments,
knowledge-soundness extractor,
adaptive ZK simulator,
authenticated fold-oracle openings,
and original KBFold QROM
analysis must all be revisited
with quantum adversaries.
\end{remark}

\begin{remark}[The honest-verifier result is substantive but bounded]
Chapter~10 constructs the concrete
commitment and affine proof at the
binary-coordinate level. The
perfect honest-sigma-verifier
simulation theorem is a real
transcript result for this
restricted model, not just an
ideal-functionality restatement.
It must not be labeled full
zero knowledge without a
malicious-verifier compiler
and the external KBFold PCS
commitment/opening simulator.
\end{remark}
